\section{Hiperparámetros y uso práctico}

\begin{frame}{Hiperparámetros principales en \texttt{sklearn.svm.SVC}}
\scriptsize
\renewcommand{\arraystretch}{1.40}
\setlength{\tabcolsep}{4.6pt}
\begin{center}
\begin{tabular}{p{2.05cm}p{4.0cm}p{6.15cm}}
\toprule
\rowcolor{ccdenavy}
\textcolor{white}{\textbf{Parámetro}} &
\textcolor{white}{\textbf{Actúa sobre}} &
\textcolor{white}{\textbf{Qué ocurre al modificarlo}}\\
\midrule
\texttt{C} & Regularización / penalización & Mayor $C$: menos tolerancia a violaciones; menor $C$: margen más ancho y más regularización.\\
\rowcolor{ccdeice!50}
\texttt{kernel} & Geometría de la frontera & \texttt{linear}, \texttt{rbf}, \texttt{poly}, \texttt{sigmoid} o función propia.\\
\texttt{gamma} & Escala del RBF/poly/sigmoid & Alto: influencia local y mayor flexibilidad. Bajo: frontera más suave. \texttt{scale} usa una regla basada en varianza.\\
\rowcolor{ccdeice!50}
\texttt{degree} & Kernel polinómico & Grado $d$; solo importa con \texttt{poly}.\\
\texttt{coef0} & Poly / sigmoid & Término independiente $r$; altera peso relativo de términos de orden alto y bajo.\\
\rowcolor{ccdeice!50}
\texttt{class\_weight} & Costos por clase & Útil con desbalance: permite penalizar más errores de una clase.\\
\bottomrule
\end{tabular}
\end{center}
\source{Géron (2022), cap. 5; API conceptual de SVC usada en el código asociado.}
\end{frame}

% ================================================================
\begin{frame}{Escalar variables no es opcional en SVM}
\small
\begin{columns}[T,onlytextwidth]
\begin{column}{0.46\textwidth}
\begin{contrastbox}{warn}{Sin escalamiento}
\small
Si una variable está en miles y otra entre 0 y 1, las distancias y productos internos quedan dominados por la primera.
\end{contrastbox}
\vspace{0.14cm}
\begin{contrastbox}{ccdeaccent}{Con \texttt{StandardScaler}}
\small
Las variables se centran y escalan antes de estimar la SVM. Para evitar fuga de información, el escalamiento debe vivir dentro de un \texttt{Pipeline}.
\end{contrastbox}
\end{column}
\begin{column}{0.50\textwidth}
\centering
\begin{tikzpicture}[scale=0.88]
\node[draw,rounded corners,fill=white,minimum width=3.2cm,minimum height=0.75cm] (raw) at (0,2.8) {datos crudos};
\node[draw,rounded corners,fill=ccdeice,minimum width=3.2cm,minimum height=0.75cm] (scale) at (0,1.45) {\texttt{StandardScaler}};
\node[draw,rounded corners,fill=ccdelight!22,minimum width=3.2cm,minimum height=0.75cm] (svc) at (0,0.1) {\texttt{SVC}};
\node[draw,rounded corners,fill=white,minimum width=3.2cm,minimum height=0.75cm] (pred) at (0,-1.25) {predicción};
\foreach \a/\b in {raw/scale,scale/svc,svc/pred}{\draw[->,thick,ccdeaccent](\a)--(\b);}
\node[text=ccdemid,font=\scriptsize,align=center] at (4.15,1.45){fit del scaler\\solo con training};
\draw[->,ccdemid] (3.3,1.45)--(1.7,1.45);
\end{tikzpicture}
\end{column}
\end{columns}
\source{Géron (2022), Figura 5-2 y ejemplo de pipeline con \texttt{StandardScaler}.}
\end{frame}

% ================================================================
\begin{frame}{Cómo ajustar $C$ y $\gamma$ sin engañarnos}
\small
\begin{columns}[T,onlytextwidth]
\begin{column}{0.49\textwidth}
\begin{contrastbox}{ccdeblue}{Flujo recomendado}
\small
\begin{enumerate}
  \item Separar train y test.
  \item Pipeline: escalamiento + SVC.
  \item Validación cruzada dentro de train.
  \item Explorar escala logarítmica para $C$ y $\gamma$.
  \item Elegir por una métrica coherente con el problema.
  \item Evaluar una sola vez en test.
\end{enumerate}
\end{contrastbox}
\end{column}
\begin{column}{0.47\textwidth}
\begin{keybox}
\textbf{Ejemplo de grilla}
\[
C\in\{0.1,1,10,100\}
\]
\[
\gamma\in\{0.1,1,10,\texttt{scale}\}
\]
\end{keybox}
\vspace{0.14cm}
\begin{ccdebox}
\small
En clasificación desbalanceada, accuracy puede ser pobre criterio. Considerar ROC-AUC, PR-AUC, recall, precision y costos de error.
\end{ccdebox}
\end{column}
\end{columns}
\source{Géron (2022), selección de modelos y cap. 5; James et al. (2013), sec. 9.6.2.}
\end{frame}

% ================================================================
\begin{frame}{Ventajas y desventajas}
\footnotesize
\begin{columns}[T,onlytextwidth]
\begin{column}{0.47\textwidth}
\begin{contrastbox}{ccdeaccent}{Ventajas}
\footnotesize
\begin{itemize}
  \item Excelente capacidad predictiva en datos pequeños/medianos.
  \item Muy útil en alta dimensión, incluso con $p\gg n$.
  \item Kernels capturan fronteras no lineales complejas.
  \item Problema convexo: solución global.
  \item Solo support vectors fijan la frontera.
  \item Regularización explícita mediante $C$.
\end{itemize}
\end{contrastbox}
\end{column}
\begin{column}{0.47\textwidth}
\begin{contrastbox}{warn}{Desventajas}
\footnotesize
\begin{itemize}
  \item No escala bien a conjuntos muy grandes, especialmente SVC kernelizado.
  \item Sensible al escalamiento de variables.
  \item $C$ y $\gamma$ pueden exigir búsqueda cuidadosa.
  \item Menos interpretable que modelos lineales o árboles simples.
  \item Predicciones probabilísticas no son nativas.
  \item Elegir un kernel inadecuado puede aumentar mucho sesgo o varianza.
\end{itemize}
\end{contrastbox}
\end{column}
\end{columns}
\source{Géron (2022), introducción del cap. 5; James et al. (2013), cap. 9.}
\end{frame}

% ================================================================
\begin{frame}{¿Cuándo elegir una SVM?}
\footnotesize
\begin{columns}[T,onlytextwidth]
\begin{column}{0.31\textwidth}
\begin{ccdebox}
\centering
\textbf{Alta dimensión}\\[0.12cm]
Texto, genómica, muchas variables respecto al número de observaciones. Un kernel lineal puede ser suficiente.
\end{ccdebox}
\end{column}
\begin{column}{0.31\textwidth}
\begin{ccdebox}
\centering
\textbf{Frontera no lineal}\\[0.12cm]
Cuando un modelo lineal falla y el tamaño muestral permite trabajar con RBF o polinómico.
\end{ccdebox}
\end{column}
\begin{column}{0.31\textwidth}
\begin{ccdebox}
\centering
\textbf{Precisión > explicación}\\[0.12cm]
Cuando importa más una frontera robusta que una interpretación directa de coeficientes.
\end{ccdebox}
\end{column}
\end{columns}
\vspace{0.22cm}
\begin{keybox}
\textbf{No es primera opción} para millones de observaciones con kernel RBF. En ese caso conviene comparar con modelos lineales escalables, árboles/boosting o aproximaciones al kernel.
\end{keybox}
\source{Géron (2022), cap. 5; James et al. (2013), aplicación de expresión génica en sec. 9.6.5.}
\end{frame}

% ================================================================
